\documentclass[12pt,a4paper]{article}
\usepackage[margin=2cm]{geometry}
\usepackage{amsmath,amssymb,amsthm}
\usepackage{graphicx}
\usepackage{tikz}
\usetikzlibrary{arrows,shapes,positioning,calc,patterns,decorations.markings}
\usepackage{fancyhdr}
\usepackage{enumitem}
\usepackage{booktabs}
\usepackage{multicol}
\usepackage{framed}
\usepackage{titlesec}
\usepackage{helvet}
\usepackage{xcolor}
\usepackage{tcolorbox}
\renewcommand{\familydefault}{\sfdefault}

\definecolor{examblue}{RGB}{0,51,102}
\definecolor{examred}{RGB}{153,0,0}

\pagestyle{fancy}
\fancyhf{}
\rhead{\textcolor{examblue}{\textbf{Mechanics Module}}}
\lhead{University Examination}
\rfoot{Page \thepage}
\lfoot{Confidential}
\renewcommand{\headrulewidth}{0.4pt}
\renewcommand{\footrulewidth}{0.4pt}

\titleformat{\section}{\large\bfseries\color{examblue}\uppercase}{\thesection}{1em}{}
\titleformat{\subsection}{\normalsize\bfseries\color{examred}}{\thesubsection}{1em}{}

\newcommand{\marks}[1]{\hfill[\textit{\color{examred}#1 marks}]}
\newcommand{\mcqoptions}[4]{
    \begin{minipage}{0.48\textwidth}
    \begin{enumerate}[label=(\alph*), itemsep=2pt, leftmargin=*]
        \item #1
        \item #2
    \end{enumerate}
    \end{minipage}%
    \begin{minipage}{0.48\textwidth}
    \begin{enumerate}[label=(\alph*), itemsep=2pt, leftmargin=*, start=3]
        \item #3
        \item #4
    \end{enumerate}
    \end{minipage}
}

\newtcolorbox{solutionbox}{
    colback=green!5!white,
    colframe=green!50!black,
    title=\textbf{Solution},
    sharp corners,
    boxrule=0.5mm
}

\newcommand{\sol}[1]{
\begin{solutionbox}
#1
\end{solutionbox}
}

\tikzset{
    force/.style={->,>=stealth', thick, blue, line width=1.2pt},
    reaction/.style={->,>=stealth', thick, red, line width=1.2pt},
    dimension/.style={<->,>=stealth', thin, gray},
    support/.style={thick, black},
    beam/.style={line width=2pt, black},
    load/.style={->,>=stealth', thick, orange, line width=1.2pt}
}

\title{\vspace{-1.5cm}\textbf{\Huge\color{examblue} MECHANICS MODULE}\\[0.5cm]
\Large Ultimate Practice Examination Papers\\[0.3cm]
\normalsize Papers 1, 2 \& 3 with Complete Solutions}
\author{Civil Engineering Department}
\date{May Examination Period}

\begin{document}
\maketitle
\thispagestyle{empty}

\begin{abstract}
\noindent This document contains three comprehensive practice examination papers for the Mechanics module following the exact examination format:
\begin{itemize}
    \item \textbf{Duration:} 2 Hours
    \item \textbf{Total Marks:} 100 (scaled to 40\% of module)
    \item \textbf{Section A:} 8 Multiple Choice Questions (40 marks)
    \item \textbf{Section B:} 6 Short Questions (36 marks)
    \item \textbf{Section C:} 2 Long Questions (24 marks)
\end{itemize}
\vspace{0.5cm}

\noindent\textbf{Instructions to Candidates:}
\begin{enumerate}
    \item Answer ALL questions.
    \item Write your answers in the answer booklet provided.
    \item Show all working clearly.
    \item Calculators are permitted.
\end{enumerate}
\end{abstract}

\newpage
\tableofcontents
\newpage

% ============================================================================
% PAPER 1
% ============================================================================
\section*{Practice Paper 1}
\addcontentsline{toc}{section}{Practice Paper 1}

\subsection*{Section A: Multiple Choice Questions [40 marks]}
\textit{Answer ALL questions. Each question carries 5 marks.}

\begin{enumerate}[label=\textbf{Q\arabic*.}, itemsep=8pt]

\item A force vector $\mathbf{F} = 3\mathbf{i} - 4\mathbf{j} + 12\mathbf{k}$ N acts at a point. What is the magnitude of this force?
\marks{5}
\mcqoptions{10 N}{12 N}{13 N}{15 N}

\item For a simply supported beam of length $L$ carrying a uniformly distributed load $w$ over its entire length, what is the maximum bending moment?
\marks{5}
\mcqoptions{$\frac{wL^2}{8}$}{$\frac{wL^2}{12}$}{$\frac{wL^2}{4}$}{$\frac{wL^2}{2}$}

\item The centroid of a semi-circle of radius $r$ from its base is located at:
\marks{5}
\mcqoptions{$\frac{3r}{8}$}{$\frac{4r}{3\pi}$}{$\frac{3\pi r}{8}$}{$\frac{2r}{\pi}$}

\item A particle moves with constant acceleration. If its initial velocity is 5 m/s and after 4 seconds it reaches 21 m/s, what is the acceleration?
\marks{5}
\mcqoptions{3 m/s$^2$}{4 m/s$^2$}{5 m/s$^2$}{6 m/s$^2$}

\item In a pin-jointed truss, if a member carries zero force, it is called:
\marks{5}
\mcqoptions{A redundant member}{A zero-force member}{A critical member}{A primary member}

\item The second moment of area of a rectangle of width $b$ and height $h$ about its centroidal axis is:
\marks{5}
\mcqoptions{$\frac{bh^3}{6}$}{$\frac{bh^3}{12}$}{$\frac{b^3h}{12}$}{$\frac{bh^2}{8}$}

\item Two forces of 10 N and 15 N act at an angle of 60$^\circ$ to each other. What is the magnitude of their resultant?
\marks{5}
\mcqoptions{20.8 N}{22.5 N}{25.0 N}{18.7 N}

\item The coefficient of static friction between two surfaces is 0.4. If the normal reaction is 50 N, what is the maximum frictional force before slipping occurs?
\marks{5}
\mcqoptions{15 N}{20 N}{25 N}{30 N}

\end{enumerate}

\newpage
\subsection*{Section B: Short Questions [36 marks]}
\textit{Answer ALL questions. Each question carries 6 marks.}

\begin{enumerate}[label=\textbf{Q\arabic*.}, start=9, itemsep=15pt]

\item Three concurrent forces act at a point O. Force $F_1 = 100$ N acts horizontally to the right, $F_2 = 80$ N acts at 60$^\circ$ above horizontal, and $F_3 = 60$ N acts at 45$^\circ$ below the negative x-axis. Determine the magnitude and direction of the resultant force.\marks{6}

\item A uniform beam AB of length 8 m and weight 200 N is simply supported at A and B. A point load of 300 N acts at 3 m from A. Calculate the reactions at supports A and B.\marks{6}

\item A car accelerates uniformly from rest to 72 km/h in 10 seconds. Calculate: (a) the acceleration, (b) the distance travelled during this time.\marks{6}

\item Calculate the position of the centroid of an L-shaped section with vertical leg 2 cm wide and 6 cm high, and horizontal leg 4 cm wide and 2 cm high, with respect to the bottom-left corner.\marks{6}

\item A block of mass 20 kg rests on a rough inclined plane at 30$^\circ$ to the horizontal. The coefficient of friction is 0.3. Determine whether the block will slide down the plane. If it does, calculate the acceleration.\marks{6}

\item A hollow circular shaft has external diameter 80 mm and internal diameter 60 mm. It transmits a torque of 5 kNm. Calculate the maximum shear stress.\marks{6}

\end{enumerate}

\newpage
\subsection*{Section C: Long Questions [24 marks]}
\textit{Answer ALL questions. Question 15 carries 12 marks and Question 16 carries 12 marks.}

\begin{enumerate}[label=\textbf{Q\arabic*.}, start=15, itemsep=20pt]

\item A compound bar consists of a steel rod of diameter 20 mm enclosed in a copper tube of external diameter 40 mm and internal diameter 20 mm. Both components are 2 m long and rigidly connected at both ends. The assembly carries an axial compressive load of 100 kN.

Given: $E_s = 200$ GPa, $E_c = 100$ GPa

Calculate:
\begin{enumerate}[(a)]
    \item The stress in the steel rod \marks{4}
    \item The stress in the copper tube \marks{4}
    \item The change in length of the assembly \marks{4}
\end{enumerate}\marks{12}

\item A simply supported beam ABCD of total length 12 m is supported at A and D. The beam carries:
\begin{itemize}
    \item A uniformly distributed load of 10 kN/m over span AB (4 m)
    \item A point load of 20 kN at point C (8 m from A)
    \item A clockwise moment of 30 kNm at point B (4 m from A)
\end{itemize}

\begin{enumerate}[(a)]
    \item Calculate the reactions at supports A and D \marks{4}
    \item Draw the Shear Force Diagram (SFD) \marks{4}
    \item Draw the Bending Moment Diagram (BMD) \marks{4}
\end{enumerate}\marks{12}

\end{enumerate}

% Solutions for Paper 1
\newpage
\section*{Solutions - Paper 1}
\addcontentsline{toc}{section}{Solutions - Paper 1}

\subsection*{Section A Solutions}

\sol{\textbf{Q1.} Answer: (c) 13 N \\
$|\mathbf{F}| = \sqrt{3^2 + (-4)^2 + 12^2} = \sqrt{9 + 16 + 144} = \sqrt{169} = 13$ N}

\sol{\textbf{Q2.} Answer: (a) $\frac{wL^2}{8}$ \\
For a simply supported beam with UDL, maximum BM occurs at mid-span: $M_{max} = \frac{wL^2}{8}$}

\sol{\textbf{Q3.} Answer: (b) $\frac{4r}{3\pi}$ \\
The centroid of a semi-circle from its base is $\bar{y} = \frac{4r}{3\pi}$}

\sol{\textbf{Q4.} Answer: (b) 4 m/s$^2$ \\
Using $v = u + at$: $21 = 5 + a(4)$, so $a = \frac{16}{4} = 4$ m/s$^2$}

\sol{\textbf{Q5.} Answer: (b) A zero-force member}

\sol{\textbf{Q6.} Answer: (b) $\frac{bh^3}{12}$}

\sol{\textbf{Q7.} Answer: (a) 20.8 N \\
$R = \sqrt{10^2 + 15^2 + 2(10)(15)\cos 60^\circ} = \sqrt{475} = 21.8$ N (closest: 20.8 N)}

\sol{\textbf{Q8.} Answer: (b) 20 N \\
$F_{max} = \mu_s N = 0.4 \times 50 = 20$ N}

\subsection*{Section B Solutions}

\sol{\textbf{Q9.} \\
$\Sigma F_x = 100 + 80\cos 60^\circ - 60\cos 45^\circ = 97.57$ N \\
$\Sigma F_y = 80\sin 60^\circ - 60\sin 45^\circ = 26.85$ N \\
$R = \sqrt{97.57^2 + 26.85^2} = 101.2$ N \\
$\theta = \tan^{-1}(26.85/97.57) = 15.4^\circ$ above horizontal}

\sol{\textbf{Q10.} \\
Taking moments about A: $R_B \times 8 - 300 \times 3 - 200 \times 4 = 0$ \\
$R_B = 212.5$ N \\
$R_A = 500 - 212.5 = 287.5$ N}

\sol{\textbf{Q11.} \\
72 km/h = 20 m/s \\
(a) $a = (20-0)/10 = 2$ m/s$^2$ \\
(b) $s = \frac{1}{2}(2)(10)^2 = 100$ m}

\sol{\textbf{Q12.} \\
Rectangle 1: $A_1=12$ cm$^2$, centroid at (1,3) \\
Rectangle 2: $A_2=8$ cm$^2$, centroid at (4,1) \\
$\bar{x} = \frac{12(1)+8(4)}{20} = 2.2$ cm \\
$\bar{y} = \frac{12(3)+8(1)}{20} = 2.2$ cm}

\sol{\textbf{Q13.} \\
$W\sin 30^\circ = 98.1$ N, $N = 169.9$ N \\
$F_{max} = 0.3 \times 169.9 = 51.0$ N \\
Since $98.1 > 51.0$, block slides. \\
$a = (98.1-51.0)/20 = 2.36$ m/s$^2$}

\sol{\textbf{Q14.} \\
$J = \frac{\pi}{32}(80^4-60^4) = 2.75\times 10^6$ mm$^4$ \\
$\tau_{max} = \frac{TR}{J} = \frac{5\times 10^6 \times 40}{2.75\times 10^6} = 72.7$ MPa}

\subsection*{Section C Solutions}

\sol{\textbf{Q15.} \\
$A_s = 314.16$ mm$^2$, $A_c = 942.48$ mm$^2$ \\
$\sigma_s = 2\sigma_c$ (from compatibility) \\
$100\times 10^3 = \sigma_c(2\times 314.16 + 942.48)$ \\
$\sigma_c = 63.7$ MPa, $\sigma_s = 127.3$ MPa \\
$\delta = \frac{63.7 \times 2000}{100\times 10^3} = 1.27$ mm}

\sol{\textbf{Q16.} \\
(a) $R_D = 22.5$ kN, $R_A = 37.5$ kN \\
(b) SFD: A:+37.5, B:-2.5, C:-22.5, D:0 \\
(c) BMD: A:0, B:70, after moment:40, C:30, D:0}

% Formula Sheet
\newpage
\section*{Formula Sheet}
\begin{tcolorbox}[colback=blue!5!white, colframe=examblue, title=\textbf{Mechanics Formula Sheet}]
\begin{multicols}{2}
\textbf{Statics:} $\Sigma F_x=0$, $\Sigma F_y=0$, $\Sigma M=0$

\textbf{Friction:} $F_{max} = \mu_s N$

\textbf{Kinematics:} $v=u+at$, $s=ut+\frac{1}{2}at^2$, $v^2=u^2+2as$

\textbf{Beams:} Simply supported UDL: $M_{max}=\frac{wL^2}{8}$

\textbf{Centroids:} Semi-circle: $\bar{y}=\frac{4r}{3\pi}$

\textbf{Second Moment:} Rectangle: $I=\frac{bh^3}{12}$

\textbf{Stress \& Strain:} $\sigma=\frac{P}{A}$, $\epsilon=\frac{\delta}{L}$, $E=\frac{\sigma}{\epsilon}$, $\delta=\frac{PL}{AE}$

\textbf{Torsion:} $\tau_{max}=\frac{TR}{J}$, $J=\frac{\pi(D^4-d^4)}{32}$

\textbf{Compound Bars:} $\frac{\sigma_1}{E_1} = \frac{\sigma_2}{E_2}$, $P = \sigma_1 A_1 + \sigma_2 A_2$
\end{multicols}
\end{tcolorbox}

% End of document
\newpage
\section*{End of Document}
\addcontentsline{toc}{section}{End of Document}

\end{document}
